\subsection{The open questions}\label{s:declaration}

A verifier who adopts the signature check must still settle four questions
raised by \citet{scher2025} (\cref{s:intro}). Does the signature stay stable
over time? How often will the check raise an alarm? How finely must workloads
be defined? And is the signature specific enough to separate declared work from
hidden work? We answer each in turn from the corpus of \cref{s:corpus}.

\paragraph{Stability over time.} The signature of decode stayed stable over the five days of recording. We recorded
decode in \DecodeRuns{} runs over \DecodeSessions{} sessions, each on a different day.
Three runs are in the reference of $s_R$, two in its fit and one in its threshold set.
Across all \DecodeRuns{} runs, $s_R$ flags
\DecodeFlagged{} of \DecodeWindows{} windows at $\tau_\alpha$, and each flagged window
is the first window of its run. On the two test-day runs, which are outside the reference, its false-positive rate is
\FprRoneDecode\%. Idle was less stable. $s_R$ flags \RoneCalIdle\% of calibration
idle, but only \FprRoneIdle\% of test idle, recorded one day earlier. Because
calibration idle scores above decode, it sets $\tau_\alpha$, and the false-positive rate on test
decode stays well below $\alpha = \AlphaPct\%$.

A reference recorded on a single day does not carry over to other days. In an exploratory check, we
fitted a reference to \SessRefWindows{} windows spread over about four minutes of one
decode run, and repeated this for each of \SessRuns{} runs, one per day. Averaged over
the three runs, this reference flags \SessOther\% of decode from the other days,
against \SessSame\% of the rest of its own run. The two runs recorded to one protocol,
two days apart, give \SessPair\% each way. This reference is not directly comparable to the reference of $s_R$ because it is shorter and flags by a leave-one-out rule rather than $\tau_\alpha$. Nevertheless, this check suggests that the reference should be recorded over several sessions, as the reference of $s_R$ is.

\paragraph{Alarm frequency.} The alarm rate sets how often inspectors must investigate.
On the test day, declared inference is decode only, and the
monitors (\cref{s:monitor}) watch $E = \ExpDecodeH$~h of it (\cref{t:t2}):
\begin{itemize}
  \item $s_R$ raises no alarm: $F = \FRoneDecode$~h$^{-1}$ (95\% upper bound
        \FupRoneDecode~h$^{-1}$), $A = \ARoneDecode$;
  \item the classifier, $s_C^{\mathrm{dev}}$, raises $F = \FClfDecode$~h$^{-1}$ (upper bound
        \FupClfDecode~h$^{-1}$), $A = \AClfDecode$.
\end{itemize}
This exposure is short, so these rates bound the alarm rate only loosely. On idle
($E = \ExpIdleH$~h), both monitors stay near zero.

\begin{table}[tp]
  \centering
  \footnotesize
  % generated by scripts/table_t2.py from results/monitor/false_alarms.json (repo_commit 1fbe999)
% 3-of-5 monitor on 1 s-hop test windows; upper bound: exact Poisson model;
% A interval: approximate onset-block bootstrap (under-covers)
\setlength{\tabcolsep}{3pt}
\begin{tabular}{@{}lrrrr@{}}
\toprule
Detector & onsets & $F$ (/h) & upper (/h) & $A$ (\%) \\
\midrule
\multicolumn{5}{@{}l}{\emph{declared decode}, $E = 0.475$ h} \\
\quad $s_R$ & 0 & 0.0 & 6.3 & 0.0 [0.0, 0.0] \\
\quad $s_C^{\mathrm{dev}}$ & 3 & 6.3 & 16.3 & 1.0 [0.0, 2.0] \\
\midrule
\multicolumn{5}{@{}l}{\emph{idle}, $E = 8.350$ h} \\
\quad $s_R$ & 1 & 0.1 & 0.6 & 0.0 [0.0, 0.1] \\
\quad $s_C^{\mathrm{dev}}$ & 0 & 0.0 & 0.4 & 0.0 [0.0, 0.0] \\
\midrule
\multicolumn{5}{@{}l}{\emph{non-AI busy}, $E = 0.950$ h} \\
\quad $s_R$ & 4 & 4.2 & 9.6 & 99.8 [99.6, 99.9] \\
\quad $s_C^{\mathrm{dev}}$ & 0 & 0.0 & 3.2 & 0.0 [0.0, 0.0] \\
\bottomrule
\end{tabular}

  \caption{False alarms on the test day, per background, with exposure $E$.
    ``upper'': the 95\% upper bound on $F$, which assumes Poisson onsets. Overlapping
    windows make onsets correlated, so this bound is approximate. The intervals on $A$
    come from a block bootstrap (\cref{a:impl}). The test day holds no declared prefill. Non-AI work is undeclared, so $s_R$ is correct to alarm
    on it.}
  \label{t:t2}
\end{table}

\paragraph{Granularity.} The reference must hold each declared workload. When declared prefill
is left out of the reference of $s_R$, $s_R$ flags \RonePrefillHoldout\% of it (one
development run). This is expected, because prefill is limited by compute and decode
by memory bandwidth, so the two load the GPU differently. With prefill
in the fit, $s_R$ flags \PrefillInFitFlagged{} of \PrefillInFitN{} prefill windows.
This rate is in-sample, because neither the calibration day nor the test day holds
prefill. One detector, fitted to decode, prefill and idle together, can therefore
accept all three. The verifier must record every declared workload, but does not
need a separate detector for each. A workload that is in the
threshold set but not in the fit instead raises the threshold (\cref{s:monitoring}).

\paragraph{Specificity.} The signature is specific, but by a thin margin. Under the complete reference, declared
windows score from \DeclBandLo{} to \DeclBandHi{} on $s_R$, and co-location, honest
training and non-AI windows from \HidBandLo{} to \HidBandHi{} (\cref{f:f7}). The two
bands overlap only at the first window of each decode run. The hidden run that comes
closest to declared work is the unthrottled co-location run. Its lowest score,
\RoneUnpacedMin, is \RoneMargin{} above \TauRoneHi, the upper end of the 95\% interval
of $\tau_\alpha$ under a bootstrap of the two calibration runs. With two runs, that end
is the threshold that calibration idle alone would set. The signature check also
does not fully separate this run from declared decode. Its AUC is \RoneAucDecode, and
\RoneDecodeAboveMin\% of test decode windows score above the run's lowest score.
